%% =====================================================================
%%  SWOT Analysis of Every Method
%%  SCOPE: One four-quadrant table per algorithm and per architectural decision.
%% =====================================================================
\section{SWOT Analysis of Every Method}
\label{sec:swot}

Every method in this system was chosen over alternatives, and every choice
bought something at a price. This section states both. The quadrants are used
in their strict sense: \emph{strengths} and \emph{weaknesses} are properties
of the method as implemented here and are therefore verifiable in the code
and the logs; \emph{opportunities} and \emph{threats} are external and
forward-looking, and the threats in particular are what the transfer to real
hardware will test.

A weakness listed below is not an admission of carelessness. It is the part
of the design that would have to change first if the corresponding assumption
failed, and knowing which part that is has more engineering value than a
clean-looking table.

% ---------------------------------------------------------------------
\begin{swottable}{Log-odds occupancy grid}
\swotcell{swotS}{Strengths}{%
Fusion is an addition, not a product and a renormalisation, so a scan is a
sweep of \texttt{+=} over an array: fast and numerically untroubled. Free
space is represented explicitly, which is what makes frontier exploration
possible at all. Clamping bounds the confidence any cell can accumulate, so a
persistent spurious return cannot become unarguable.}
&
\swotcell{swotW}{Weaknesses}{%
The independence assumption between cells is false: a wall is a correlated
structure, and treating its cells as independent lets a thin obstacle be
eroded by beams that pass near it. Resolution is fixed at \SI{0.10}{m}, which
is coarse relative to a \SI{0.80}{m} aisle. A moving obstacle leaves a smear
that decays only as fast as the free-space evidence accumulates.}
\\ \hline
\swotcell{swotO}{Opportunities}{%
The same grid accepts a second sensor without structural change: ultrasonic
returns, which reflect off glass where lidar does not, fuse as additional
log-odds increments. Multi-resolution variants exist if the aisle geometry
ever demands them.}
&
\swotcell{swotT}{Threats}{%
Glass. A \SI{2}{D} lidar at glancing incidence on a pane returns nothing, so
the wall is not merely mis-mapped, it is absent, and the grid will confidently
report free space where there is a barrier. This is the single most dangerous
sim-to-real gap in the system, and no mapping parameter addresses it.}
\\
\end{swottable}

% ---------------------------------------------------------------------
\begin{swottable}{A* with an octile heuristic on an inflated grid}
\swotcell{swotS}{Strengths}{%
Optimal for the given grid and cost function, and complete: if it returns no
path, no path exists at that inflation. Inflating obstacles by the footprint
radius reduces the robot to a point, which is what makes an \SI{0.80}{m}
aisle tractable for a grid planner at all. The octile heuristic is admissible
on an eight-connected grid, so optimality survives.}
&
\swotcell{swotW}{Weaknesses}{%
It plans for a point, so the resulting path takes no account of the robot's
orientation --- acceptable only because the base is holonomic. Inflation is
isotropic, which over-inflates along the aisle where the base is narrow and
under-inflates across it. Replanning is from scratch; there is no incremental
repair.}
\\ \hline
\swotcell{swotO}{Opportunities}{%
D*~Lite would give incremental replanning at a fraction of the cost when the
map changes locally. A costmap gradient rather than a binary inflation would
let the path prefer the aisle centre instead of merely staying legal.}
&
\swotcell{swotT}{Threats}{%
Inflation radius and aisle width are in direct tension here: the margin that
guarantees clearance can close the only route. A greenhouse whose aisles are
narrower than this one --- and \SI{0.80}{m} is already tight --- would make
the planner report failure rather than risk, which is safe but useless.}
\\
\end{swottable}

% ---------------------------------------------------------------------
\begin{swottable}{Pure pursuit path following}
\swotcell{swotS}{Strengths}{%
One parameter of real consequence, the lookahead, with a physical meaning
that can be reasoned about rather than tuned blindly. Naturally smooth: it
cannot produce the discontinuous commands a waypoint-snapping follower does.
Cheap enough to run at \SI{20}{Hz} beside everything else.}
&
\swotcell{swotW}{Weaknesses}{%
Corner cutting is intrinsic --- the lookahead point is the target, not the
path --- and the amount cut scales with the lookahead, so the parameter
trades tracking accuracy against stability with no setting that gives both.
No feedforward: it reacts to the path, it does not anticipate it. It has no
model of the actuators, so it will happily command a twist the base cannot
produce.}
\\ \hline
\swotcell{swotO}{Opportunities}{%
An adaptive lookahead scaled by speed and by local curvature is a small
change with a measurable effect. The controller is also the natural point of
comparison for the model-predictive alternative discussed in
Sec.~\ref{sec:discussion}: identical path, identical disturbance, two
control laws.}
&
\swotcell{swotT}{Threats}{%
Wheel slip on a wet greenhouse floor breaks the assumption that the commanded
twist is the achieved twist, and pure pursuit has no mechanism to notice. In
an aisle with \SI{0.21}{m} of margin per side, uncorrected lateral drift
reaches the gutter well before the next replan.}
\\
\end{swottable}

% ---------------------------------------------------------------------
\begin{swottable}{Colour-threshold fruit detection}
\swotcell{swotS}{Strengths}{%
Runs at frame rate on a CPU with no model, no training set and no dependency
to install --- which is exactly why it exists: it made the whole mission
testable months before a learned detector could be. The excess-red form used
here is invariant to illumination \emph{scale}, so it survives shade.}
&
\swotcell{swotW}{Weaknesses}{%
\measured{69\%} recall and \measured{66 of 111} map estimates matching no
real berry. It cannot be tuned out of this: the test is invariant to
illumination scale but not to illumination \emph{colour}, and late-afternoon
light through glass lowers $G/R$ for the entire scene until every surface
passes. It has no notion of shape, texture or context, so a red pipe is
fruit.}
\\ \hline
\swotcell{swotO}{Opportunities}{%
It remains a legitimate baseline, and a strong one for the report: a learned
detector that does not beat it \emph{on the same images} has proved nothing.
It also survives as a cheap pre-filter to bound the region a heavier model
must examine.}
&
\swotcell{swotT}{Threats}{%
Real greenhouse illumination varies far more than the simulated world's, and
in the direction the threshold is weakest. Every hour of the day is a
distribution shift, and the failure is silent --- the detector reports
confident fruit, not an error.}
\\
\end{swottable}

% ---------------------------------------------------------------------
\begin{swottable}{Corroborated fruit map (two passes and a baseline)}
\swotcell{swotS}{Strengths}{%
It gives the mission the one thing no detector improvement can supply: the
ability to \emph{disbelieve} a single observation. The two conditions reject
different things --- the pass count rejects the transient, the viewpoint
baseline rejects the view-dependent artefact such as a specular highlight,
which a detector at \SI{15}{Hz} would otherwise corroborate thirty times from
effectively one place. It is deliberately ignorant of detector confidence, so
it keeps working unchanged when the threshold is replaced by a model.}
&
\swotcell{swotW}{Weaknesses}{%
It resolves clusters, not berries, and this is a hard limit rather than a
tuning choice: measured fruit position error is \SI{0.21}{m} while berries on
one plant sit \SI{0.05}{m} to \SI{0.20}{m} apart, so the error exceeds the
spacing. It also costs a scouting pass --- corroboration is not free, it is
paid for in survey time. A berry visible from only one side of the aisle can
never accumulate a baseline and is therefore never confirmed.}
\\ \hline
\swotcell{swotO}{Opportunities}{%
The same structure accepts a ripeness estimate per cluster, turning the
harvest route into a scheduling problem over time rather than a tour over
positions. Cluster persistence across sessions would let the robot return to
fruit that was not ready yesterday.}
&
\swotcell{swotT}{Threats}{%
The rule assumes the fruit does not move between passes. Wind through open
vents, or a picker working the same row, breaks that assumption in the
direction that costs recall rather than precision --- which is the safer
direction, but still a loss.}
\\
\end{swottable}

% ---------------------------------------------------------------------
\begin{swottable}{Command-level protective stop with a single owner}
\swotcell{swotS}{Strengths}{%
Enforcement is structural rather than procedural: one node writes
\topicname{cmd\_vel} and no publisher can reach the wheels without passing
through it, so the guarantee does not depend on every future contributor
remembering. Clearance is exact geometry validated against brute force
(Sec.~\ref{sec:safety}), and the same function is used by the guard, by the
escape manoeuvre and by the commissioning tool that certifies the distance.}
&
\swotcell{swotW}{Weaknesses}{%
It is a single point of failure by construction --- that is the design ---
so a defect in it is a defect in all containment. It also brakes on a
geometric criterion alone, with no notion of what the obstacle is: a leaf and
a person produce the same stop. Cancelling translation while preserving
rotation is correct for escape but means the robot can still turn in place
against something it is already touching.}
\\ \hline
\swotcell{swotO}{Opportunities}{%
The structure accepts a second, independent channel --- bumper strips or
ultrasonic rangefinders --- without redesign, which is the normal route to a
performance level a certification body will accept.}
&
\swotcell{swotT}{Threats}{%
A software-only protective stop on a kinematically driven base does not
transfer as-is to a machine with real inertia: the stopping distance becomes
a dynamic quantity, and ISO~13855 requires it to be measured, not assumed.
Commissioning stage~0 exists precisely to measure it, and every clearance
figure in this report is provisional until it does.}
\\
\end{swottable}

% ---------------------------------------------------------------------
\begin{swottable}{Discrete state feedback for visual alignment}
\swotcell{swotS}{Strengths}{%
It removes a steady-state error that no proportional gain could: replayed
against the logged stalls it reaches exactly zero offset where the previous
law left \SI{0.16}{} to \SI{0.29}{}. The design is analytic --- settling time
and damping in, gains out --- so the behaviour is specified rather than
discovered, and the closed-loop poles are asserted in the regression suite.}
&
\swotcell{swotW}{Weaknesses}{%
The model is deliberately minimal, and its gain depends on range, so the
poles are only verified stable across the \SI{0.3}{} to \SI{2.5}{m} band that
matters here. It assumes the measured offset is the true offset, which the
detector's false positives violate. It is single-input, single-output, so it
cannot coordinate the lateral and heading corrections it is implicitly
splitting.}
\\ \hline
\swotcell{swotO}{Opportunities}{%
The multivariable extension --- four wheel commands reduced to three body
degrees of freedom, with the decoupling done in the gain matrix --- is a
direct continuation and the natural next contribution.}
&
\swotcell{swotT}{Threats}{%
No control law fixes an actuator limit. At \SI{2}{m} an offset of $0.68$
requires \SI{1.36}{m} of lateral travel, which is \SI{14}{s} at the
\SI{0.10}{m/s} ceiling against a \SI{12}{s} timeout: the alignment was never
going to succeed at any gain. The answer there is to refuse the station, not
to tune the controller, and mistaking one for the other cost several runs.}
\\
\end{swottable}

% ---------------------------------------------------------------------
\begin{swottable}{Simulation-first development with a digital twin}
\swotcell{swotS}{Strengths}{%
It permits failures that would be unacceptable on hardware, and it supplies
ground truth --- the world knows where all \measured{127} berries are ---
which makes every claim in Sec.~\ref{sec:results} a measurement rather than
an impression. Iteration is limited by a twenty-minute run, not by a
greenhouse's availability or a season.}
&
\swotcell{swotW}{Weaknesses}{%
The twin is a model of what was already understood, so it cannot surface a
phenomenon nobody thought to include, and the sensor models are optimistic in
exactly the places that matter: the lidar sees glass, the camera has no
motion blur, the floor does not slip. Success in simulation is evidence about
the software, not about the robot.}
\\ \hline
\swotcell{swotO}{Opportunities}{%
The gap is narrowable in identified directions: an actuator model with
transport delay and rate limits, domain-randomised illumination, and
odometry drift injected from a calibrated model rather than assumed. The
first of these is implemented and awaiting validation.}
&
\swotcell{swotT}{Threats}{%
The characteristic failure of this method is confidence: a system that has
only ever succeeded in simulation looks finished. The commissioning stages
exist to make that confidence expensive to hold, by requiring a measured
result on hardware before each capability is declared.}
\\
\end{swottable}
